% migrate-to: market-structure-execution-quality
\subsection{Smart Order Routing}

\begin{definition} \hlt{Execution-Side SOR}\\
In execution, a smart order router chooses where each child order should go after the parent algorithm has chosen when and how much to trade. It is the venue-optimisation layer inside the execution stack.
\end{definition}

\begin{definition} \hlt{Marketable Price Across Venues}\\
With several books, marketability is venue-specific. A buy order can be marketable on one venue because its limit reaches the local ask, while not marketable elsewhere if the same instrument is offered at a higher price.
\end{definition}

\begin{taxonomy} \hlt{Routing Logic Segmentation}
\begin{enumerate}[label=\roman*.]
\setlength{\itemsep}{0pt}
\item \hlt{observable liquidity}: quotes, depth, queue position, hidden-liquidity signals, and auction state;
\item \hlt{fill probability}: venue hit rate, cancellation speed, adverse selection, and reject history;
\item \hlt{net economics}: explicit fees, maker-taker rebates, clearing, and broker-specific schedules;
\item \hlt{latency}: feed delay, gateway delay, matching-engine delay, and stale-quote risk;
\item \hlt{information cost}: signalling from sweeping, probing, or repeatedly visiting the same venue.
\end{enumerate}
\end{taxonomy}

\begin{process} \hlt{Child-Order Routing Loop}
\begin{enumerate}[label=\roman*.]
\setlength{\itemsep}{0pt}
\item receive the child quantity, side, limit, and urgency from the execution algorithm;
\item score venues by expected net cost and completion probability;
\item split or sequence messages across venues;
\item update venue scores from fills, partial fills, rejects, and missed liquidity.
\end{enumerate}
\end{process}

\begin{principle} \hlt{Best Execution Differs by Regime}\\
US routing leans on protected consolidated quotes. European brokers choose the consolidation layer, so routing quality includes venue access, data quality, and auditability.
\end{principle}

\begin{caveat} \hlt{Displayed Best Price Is Not Enough}\\
A venue with the best quote can still be a poor destination if the queue is unstable, the quote is stale, fees are high, or fills are followed by adverse price moves.
\end{caveat}

\begin{mechanism} \hlt{Active Sweep Segmentation}
\begin{enumerate}[label=\roman*.]
\setlength{\itemsep}{0pt}
\item \hlt{local best price}: cheapest visible venue first;
\item \hlt{consolidated price}: split across venues to minimise average price;
\item \hlt{policy coverage}: restrict destinations to venues covered by the broker mandate;
\item \hlt{residual choice}: when venues tie at a level, choose by fee, latency, and fill reliability.
\end{enumerate}
\end{mechanism}

\begin{caveat} \hlt{Fee-Adjusted Price Priority}\\
Let $P$ be the price on $V$, $\delta$ the relative price increase on $V'$, $Q$ the size, $f$ the fee on $V$, and $\phi$ the fee saving on $V'$. For a buy order, $V'$ is cheaper when $P(1+\delta)(1+f-\phi)Q < P(1+f)Q$.\\
Roughly, the price penalty must be below the fee saving.
\end{caveat}

\begin{caveat} \hlt{Order Preferencing}\\
A broker may prefer a dealer, affiliate, or venue because of service, economics, or reciprocal flow rather than the displayed quote alone. Best-execution review must therefore compare all-in fills, not just routing intent.
\end{caveat}

\begin{proposition} \hlt{Payment-for-Order-Flow Routing Conflict}\\
Let $C_v$ be expected client execution cost on venue or dealer $v$, and let $h_v$ be broker revenue from routing there. Client-optimal routing chooses $\arg\min_v C_v$; broker-private routing chooses $\arg\min_v(C_v-h_v)$. A conflict exists whenever some $a,b$ satisfy $C_a<C_b$ but $C_a-h_a>C_b-h_b$.
\begin{proof}
The client pays execution cost $C_v$.\\
The broker internalises both client execution quality and routing revenue, so its private objective subtracts $h_v$.\\
The displayed inequalities state that $a$ is better for the client while $b$ is better for the broker.
\end{proof}
\end{proposition}

\begin{caveat} \hlt{Latency Makes Liquidity Probabilistic}\\
Displayed depth is probabilistic: stale or duplicated quotes can vanish before the router arrives.
\end{caveat}

\begin{principle} \hlt{Liquidity Seeking Objective}\\
A liquidity seeker scores the quantity likely to be available on arrival, using visible depth, hidden-liquidity signals, latency, duplicate quotes, fees, and fill rates.
\end{principle}

\begin{taxonomy} \hlt{Passive Posting Segmentation}
\begin{enumerate}[label=\roman*.]
\setlength{\itemsep}{0pt}
\item \hlt{queue rank}: post more where fewer shares stand ahead;
\item \hlt{distance to touch}: prefer venues closer to becoming marketable;
\item \hlt{trading rate}: prefer venues where opposing flow consumes the level faster;
\item \hlt{resting-time target}: split to equalise expected completion time.
\end{enumerate}
\end{taxonomy}

\begin{definition} \hlt{Minimum Resting-Time Split}\\
Let $D$ be total passive child size, $Q_n$ queue ahead, $r_n$ opposing-flow rate, $f_n$ posted fraction, and $\tau^\star$ common completion time. For posted size $x$ on venue $n$, $\tau_n(x)=\frac{x+Q_n}{r_n}$ is the expected completion time, so equal-time posting gives $\tau^\star=\frac{D+\sum_n Q_n}{\sum_n r_n}$ and $f_n=\frac{\tau^\star r_n-Q_n}{D}$. If $f_n<0$, post nothing on venue $n$.
\end{definition}

\begin{table}[htbp]
\centering
\scriptsize
\stepcounter{table}
\renewcommand{\arraystretch}{1.12}
\setlength{\tabcolsep}{4pt}
\begin{tabularx}{\textwidth}{@{}>{\RaggedRight\arraybackslash}p{0.28\textwidth} >{\RaggedRight\arraybackslash}p{0.18\textwidth} >{\RaggedRight\arraybackslash}p{0.18\textwidth} >{\RaggedRight\arraybackslash}X@{}}
\hline
\textbf{Criterion} & \textbf{1st venue} & \textbf{2nd venue} & \textbf{3rd venue} \\
\hline
Rank in queue & better queue & next queue & deepest queue \\
Distance to best & nearest touch & next touch & farthest touch \\
Trading rate & fastest flow & middle flow & slowest flow \\
Minimum resting time & highest adjusted rate & next adjusted rate & residual venue \\
\hline
\end{tabularx}

\vspace{0.25em}
{\footnotesize Table \thetable. Passive routing criteria can rank the same venues differently.}
\end{table}

\begin{process} \hlt{Liquidity-Seeker Build}
\begin{enumerate}[label=\roman*.]
\setlength{\itemsep}{0pt}
\item estimate hidden liquidity, mirrored orders, latency, fees, and adverse selection;
\item remove liquidity when the probability-adjusted net price is attractive;
\item post passively where queue, distance, and trading-rate criteria support fast completion;
\item keep the parent order inside the execution envelope set by the benchmark layer.
\end{enumerate}
\end{process}
